\lysh{19038_SOLUTION}{1}
\textbf{ΛΥΣΗ}

\begin{enumerate}[label=\alph*)]
    \item Για τη γωνία $\omega$ που σχηματίζει το διάνυσμα $\vec{\beta}$ με τον άξονα $x'x$ είναι:
    \[ \text{εφ}\omega = \frac{y_{\vec{\beta}}}{x_{\vec{\beta}}} = \frac{1}{-1} = -1 \]
    Αφού είναι $x_{\vec{\beta}} < 0$ και $y_{\vec{\beta}} > 0$, για τη γωνία $\omega$ θα ισχύει:
    \[ \frac{\pi}{2} < \omega < \pi \]
    Άρα, $\omega = \pi - \frac{\pi}{4} = \frac{3\pi}{4}$.

    \item Είναι:
    \[ |\vec{\gamma}| = \sqrt{x_{\vec{\gamma}}^2 + y_{\vec{\gamma}}^2} = \sqrt{50} = 5\sqrt{2} \]
    \[ |\vec{\beta}| = \sqrt{x_{\vec{\beta}}^2 + y_{\vec{\beta}}^2} = \sqrt{2} = \sqrt{2} \]
    Επομένως, $|\vec{\gamma}| = 5|\vec{\beta}|$.

    \item Ζητάμε τους πραγματικούς αριθμούς $\lambda, \mu$ για τους οποίους ισχύει η ισότητα:
    \[ \vec{\gamma} = \lambda\vec{\alpha} + \mu\vec{\beta} \]
    Έχουμε διαδοχικά:
    \[ (-5,-5) = \lambda(2,3) + \mu(-1,1) \]
    \[ (-5,-5) = (2\lambda-\mu, 3\lambda + \mu) \]
    Επομένως, είναι:
    \[ 2\lambda-\mu = -5 \quad (1) \text{ και } 3\lambda + \mu = -5 \quad (2) \]
    Με πρόσθεση κατά μέλη προκύπτει ότι:
    \[ 5\lambda = -10 \text{ ή } \lambda = -2 \]
    Αντικαθιστούμε στην (2), οπότε:
    \[ 3(-2) + \mu = -5 \text{ ή } \mu = -5 + 6 = 1 \]
    Άρα, το διάνυσμα $\vec{\gamma}$ γράφεται ως εξής:
    \[ \vec{\gamma} = -2\vec{\alpha} + \vec{\beta} \]
\end{enumerate}